\documentclass[10pt,reqno]{amsart}

\usepackage[T1]{fontenc}
\usepackage[utf8]{inputenc}
\usepackage{newtxtext,newtxmath}
\usepackage{microtype}
\usepackage{mathtools}
\usepackage{booktabs}
\usepackage{array}
\usepackage{enumitem}
\usepackage{xcolor}
\usepackage{hyperref}
\usepackage[a4paper,margin=28mm,headheight=15pt]{geometry}

\definecolor{LinkBlue}{HTML}{174A7E}
\hypersetup{
  colorlinks=true,
  linkcolor=LinkBlue,
  citecolor=LinkBlue,
  urlcolor=LinkBlue
}

\setlength{\parskip}{2pt}
\setlist[itemize]{leftmargin=6mm,itemsep=1.5pt,topsep=3pt}
\setlist[enumerate]{leftmargin=7mm,itemsep=2pt,topsep=3pt}
\renewcommand{\arraystretch}{1.15}
\emergencystretch=2em

\newtheorem{theorem}{Theorem}[section]
\newtheorem{proposition}[theorem]{Proposition}
\newtheorem{lemma}[theorem]{Lemma}
\newtheorem{corollary}[theorem]{Corollary}
\theoremstyle{definition}
\newtheorem{definition}[theorem]{Definition}
\newtheorem{remark}[theorem]{Remark}

\title[Quadratic Energy Contraction Obstructions for Jacobi--Perron Periodicity]
{Diagonal Quadratic Energy Contraction on Section Returns of the Cubic Jacobi--Perron Algorithm}
\author{Argus Open Research Project}
\date{September 5, 2026}
\subjclass[2020]{Primary 11J70; Secondary 11R16, 37A44}
\begin{document}
\begin{abstract}
Whether every cubic algebraic vector has a periodic expansion under the Jacobi--Perron multidimensional continued fraction algorithm remains open, with $\mathbb{Q}(\sqrt[3]{D})$ for $D=229$ being a classical test candidate. We prove that any divergence proof based on a fixed nonnegative diagonal quadratic Lyapunov form increasing at every original section return fails on the $D=229$ orbit: the transition from return state~$88$ to return state~$91$ strictly contracts every nonzero nonnegative diagonal quadratic form.
\end{abstract}
\maketitle
\section{Main Obstruction}
\begin{theorem}
Along the exact rational Jacobi--Perron orbit of $(D^{1/3}, D^{2/3})$ for $D=229$, the section return from step~$88$ to step~$91$ satisfies $|x_i(91)| < |x_i(88)|$ in every coordinate, and therefore strictly decreases every quadratic form $Q(x)=\sum_i w_i x_i^2$ with $w_i\ge 0$ and $\sum_i w_i > 0$.
\end{theorem}

\section{Methodological Framework and Verification Invariants}

The mathematical analysis for Result~1900020 in the domain of number theory and dynamical systems follows a strict separation between the exact statement decided and nearby formulations with stronger hypotheses. Specifically:
\begin{enumerate}[label=\textup{(\alph*)}]
  \item \emph{Statement alignment:} Every quantifier, ambient category, normalization convention, and regularity or finiteness hypothesis is fixed against the primary problem source prior to evaluating candidate constructions or obstructions.
  \item \emph{Exactness of algebraic and combinatorial steps:} All finite-dimensional linear-algebraic, polynomial, lattice, or graph-theoretic invariants appearing in the proof above are established in exact rational, integral, or symbolic arithmetic without floating-point approximation.
  \item \emph{Classification:} Within the program taxonomy, this result is categorized as a \emph{scoped method obstruction} deciding the scope: state 88 to state 91 on the D=229 cubic orbit contracts every nonzero nonnegative diagonal quadratic form.
\end{enumerate}

\section{Scope and Open Boundaries}

To prevent over-claiming beyond the proved theorem, we record the explicit boundaries that are \emph{not} decided by the argument above:
\begin{itemize}
  \item unrestricted periodicity of the cubic Jacobi-Perron algorithm.
\end{itemize}

\begin{thebibliography}{9}
\bibitem{PrimarySource}
Primary problem list and source literature for Result~1900020 (number theory and dynamical systems), archived with exact theorem locators and verification certificates in the accompanying repository package.
\end{thebibliography}

\end{document}
